\item Está trabajando en un laboratorio de materia condensada para su proyecto de graduación. Varios de los proyectos en curso utilizan helio líquido, que está contenido en un recipiente con aislamiento térmico que puede contener hasta un máximo de $V_{\text{máx}} = 240\text{ L}$ del líquido en $T = 4.20\text{ K}$. Debido a que ya se ha utilizado parte del helio líquido, alguien le pide que verifique si hay suficiente para el día siguiente, en el que cuatro grupos experimentales diferentes lo necesitarán. No está seguro de cómo medir la cantidad de líquido restante, por lo que inserta una varilla de aluminio de longitud $L = 2.00\text{ m}$ y con un área de sección transversal de $A = 2.50\text{ cm}^2$ en el recipiente. Al ver qué parte del extremo inferior de la varilla está helada al extraerla, puede estimar la profundidad del helio líquido. Sin embargo, después de insertar la varilla, uno de los experimentadores lo llama para realizar una tarea y se olvida de la varilla, dejándola en el helio líquido hasta la mañana siguiente. ¿Cuánto helio líquido está disponible para los experimentos del día siguiente? (El aluminio tiene una conductividad térmica de $3100\text{ W/m}\cdot\text{K}$ a 4.20 K; ignore su variación de temperatura. La densidad del helio líquido es de $125\text{ kg/m}^3$.) Suponga que el helio gaseoso puede escapar de la parte superior del recipiente.
